\section{Evaluation Blind Spots}

\subsection{Problem Definition}

For a query $q$, let $\gold(q)=\{g_1,\ldots,g_m\}$ denote human-gold
evaluation dimensions and $\pred(q)=\{r_1,\ldots,r_n\}$ denote dimensions
generated by a model. A blind spot occurs when a gold dimension $g_i$ is not
semantically covered by any generated criterion:
\[
  \max_j \cos(e(g_i), e(r_j)) < \tau.
\]
This definition is intentionally dimension-level rather than score-level. A
judge may still choose the correct answer on a particular pair while missing a
dimension that would matter on nearby examples, harder counterexamples, or
distribution shifts. Measuring blind spots therefore probes the reliability of
the evaluative basis before it is collapsed into a scalar judgment.

\subsection{Diagnostic Evidence}

We use a frozen motivation diagnostic on 100 hard-gold RubricBench examples.
The data, generation, verification, bootstrap confidence intervals, and claim
matrix pass the repository's \emph{minimal motivation} gate, and that gate
checks the exact frozen snapshot before these numbers are used in the
motivation section. This gate is separate from the full real-run C0--C14
readiness matrix: the diagnostic can motivate the problem, but it cannot make a
trained-method row paper-facing. At the
default protocol thresholds $\tau=0.75$ for coverage and $\tau'=0.85$ for
redundancy, the base single-model evaluation-criteria policy achieves only 0.3692 mean
coverage and leaves a 0.6308 mean blind-spot rate. The bootstrap 95\%
confidence interval for the blind-spot rate is 0.5785--0.6823, well above the
non-triviality gate of 0.2.
This is not driven by a small number of outliers: the median query-level blind
rate is 0.6667, 74 of 100 diagnostic queries cover at most half of their
human-gold dimensions, and 21 queries have zero covered gold dimensions. The
same run reports 0.0768 redundancy and 0.1227 hallucination/invalidity,
showing that the failure is not merely a formatting artifact.
The conclusion is also not an artifact of a single semantic threshold: in the
same sweep, the blind-spot rate remains 0.4401 under a looser coverage
threshold $\tau=0.70$ and rises to 0.7965 under a stricter threshold
$\tau=0.80$.

This diagnostic motivates the rest of the paper. The failure is not that an
evaluation-criteria policy occasionally omits a niche item, but that most hard-gold
queries lose a large fraction of the human evaluative basis before any judge
score is computed. In this setting, final preference accuracy can hide an
upstream coverage gap: the judge may be right on an easy instance while still
lacking dimensions needed for nearby or harder cases. BSC is therefore used as
a dimension-level diagnostic and reward, rather than as another surface metric
for judging criteria verbosity.

\begin{table}[t]
\centering
\begin{tabular}{lrrrr}
\toprule
Setting & Coverage$\uparrow$ & Blind$\downarrow$ & Redundancy$\downarrow$ & Hallucination$\downarrow$ \\
\midrule
Base, $\tau=0.75,\tau'=0.85$ & 0.3692 & 0.6308 & 0.0768 & 0.1227 \\
\bottomrule
\end{tabular}
\caption{Minimal hard-gold diagnostic on 100 RubricBench examples. The
central evidence is the high blind-spot rate: the model omits most
human-gold evaluation dimensions even after producing plausible criteria.}
\label{tab:minimal_blindspot}
\end{table}

\subsection{Blind Spots Are Structured}

The missed dimensions are not uniformly random. A blind-spot attribution pass
maps uncovered human-gold dimensions into seven coarse categories: factuality,
completeness, constraint following, safety, domain knowledge, evidence
grounding, and intent/reasoning. On the 100-example diagnostic split,
constraint-following dimensions are especially under-covered, with a blind
rate of 0.8226, followed by intent/reasoning at 0.7692 and completeness at
0.6219. This supports the paper's central motivation: model-generated criteria
often capture common surface dimensions but miss long-tail constraints and
implicit evaluative requirements. This attribution is descriptive motivation
evidence only; category-level reduction claims require the later C12
dimension-transition audit and trained-method gates.

\begin{table}[t]
\centering
\begin{tabular}{lrrr}
\toprule
Category & Gold Dims. & Coverage$\uparrow$ & Blind$\downarrow$ \\
\midrule
Factuality & 69 & 0.4203 & 0.5797 \\
Completeness & 283 & 0.3781 & 0.6219 \\
Constraint following & 124 & 0.1774 & 0.8226 \\
Safety & 7 & 0.7143 & 0.2857 \\
Domain knowledge & 58 & 0.4483 & 0.5517 \\
Evidence grounding & 15 & 0.6667 & 0.3333 \\
Intent/reasoning & 13 & 0.2308 & 0.7692 \\
\bottomrule
\end{tabular}
\caption{Blind-spot attribution on the minimal diagnostic split. The largest
gaps appear in constraint following and intent/reasoning, which are precisely
the dimensions that generic judge criteria tend to under-specify.}
\label{tab:blindspot_categories}
\end{table}

\subsection{Not Just a Criteria-Budget Effect}

One trivial explanation for a larger measured coverage value is that the
model simply outputs more criteria. We therefore evaluate a criteria-budget curve by
truncating the generated list to $K$ items and recomputing BSC. Under a larger
criteria budget, coverage moves from 0.2338 at $K=3$ to 0.3692 at $K=10$, but
saturates afterward because the base policy produced 842 total criteria over
the 100 examples. The saturation shows that blind spots persist even after
allowing substantially more criteria items; simply making criteria lists longer
is insufficient without a reward that targets uncovered human dimensions and
discourages redundancy.
